% ============================================================================
% sections/sec2_literature.tex
% ============================================================================
\section{Literature Survey}
\label{sec:literature}

\subsection{Introduction to Literature Survey}

This survey covers six papers (2025--2026) that, together, define a
working threat model and defense landscape for retrieval poisoning in RAG
systems: a foundational spectral defense (DRS), a query-time reshaping
defense (ShieldRAG), a two-layer adversarial-training-plus-counterfactual
defense (RAGuard/ZKIP), a compound query-and-corpus attack (PIDP-Attack),
a stealth fluency-preserving attack (SilentRetrieval), and a three-ring
layered defense with an explicitly reported negative result
(TriShieldRAG). All six were verified against the actual paper text (not
a secondary summary) before being cited in this report; exact identifiers
appear in the References section.

\subsection{Related Work}

\subsubsection{DRS --- Directional Relative Shifts (ICLR 2025, under double-blind review)}

``Understanding Data Poisoning Attacks for RAG: Insights and Algorithms''
observes that effective poisoning attacks tend to create unusually large
shifts along the directions in which the clean-data distribution has the
smallest variance. Building on that insight, the paper proposes DRS
(Directional Relative Shifts): fit a PCA subspace on a trusted reference
corpus, identify the lowest-variance directions, and flag any document
whose projection onto those directions is abnormally large. Its stated
strengths are a principled statistical foundation and stronger filtering
than perplexity-, Euclidean-distance-, or embedding-norm-based baselines
in the paper's own medical-RAG evaluation setting. Its stated limitations
are a targeted-query assumption that weakens as the query space grows,
the fact that it only screens documents at ingestion and never assesses
whether the answer eventually produced is itself trustworthy, and a
measured roughly 15\% drop in detection under an adaptively-regularised
poison.

\subsubsection{ShieldRAG --- ``Push and Pull'' (ACM Transactions on Information Systems, 2026)}

``Push and Pull: Defending against Retrieval Poisoning Attacks via
Embedding Space Reshaping'' proposes reshaping the retrieval embedding
space at query time: push the query embedding away from
minority/suspicious signals and pull it toward majority/benign signals,
using a sliding-window explanation mechanism, keyword aggregation, and
query-targeting optimisation. Its stated strength is that it requires no
modification to the corpus and operates purely at query time. Its stated
limitation, acknowledged directly in the paper, is that the whole
mechanism assumes benign information is the majority of what is
retrieved --- under collusion (two or more poison documents), the
malicious side can itself become the majority, at which point
majority-consensus reshaping reinforces the wrong answer instead of
correcting it. This project's own baseline implementation of ShieldRAG
reproduces that limitation directly: measured across every attacked case
in this evaluation, ShieldRAG-only never once diverges from plain Vanilla
RAG's answer (0 divergent outcomes), because push-pull reweighting only
reinforces whichever answer already holds the plurality and has no
external signal telling it that plurality might be wrong.

\subsubsection{RAGuard / ZKIP (arXiv:2607.26339, 2026)}

``RAGuard: A Layered Defense Framework for Retrieval-Augmented Generation
Systems Against Data Poisoning'' combines two layers: an adversarially
fine-tuned dense retriever that learns to downrank passages resembling
synthetic poison, and the Zero-Knowledge Inference Patch (ZKIP) --- a
label-free, black-box, leave-one-out filter that decodes a reference
answer with the full retrieved context, decodes again with one document
removed, and flags that document if its removal substantially shifts the
answer's semantics or reduces the model's output-entropy. ZKIP needs no
poison labels, no gold answers, and no model internals, and the paper
reports it drives measured attack success to 0.000 in every defended
configuration it tests. The paper is equally direct about ZKIP's own
limitation, naming it explicitly as the central known weakness of
single-pass ZKIP: \enquote{When two or more poisons in the top-$k$ assert
the same false fact, removing any one leaves the answer unchanged; the
LOO signal for each is muted and no document is flagged.} This project's
Ring~3 (\Cref{sec:ring3}) is a direct response to that stated gap.

\subsubsection{PIDP-Attack (arXiv:2603.25164, 2026)}

``PIDP-Attack: Combining Prompt Injection with Database Poisoning Attacks
on Retrieval-Augmented Generation Systems'' demonstrates a compound
attack that pairs a runtime query-path prompt injection with corpus-path
poisoning, so that neither surface alone needs to succeed --- the query
suffix steers retrieval toward the poisoned documents, and the poisoned
documents then carry the false claim into the generator's context. Its
significance for this project is architectural rather than defensive: it
is the direct justification for treating the query path (Ring~0) as a
first-class, independently-screened surface rather than trusting that
corpus-side defenses alone are sufficient.

\subsubsection{SilentRetrieval (ACM SIGKDD 2026 / arXiv:2605.28074)}

``SilentRetrieval: Hijacking Retrieval-Augmented Generation via
Semantically-Preserving Adversarial Data Poisoning'' introduces a
two-stage stealth construction: Coordinated Beam Search (CBS), a
multi-token joint optimisation that keeps a topically relevant host
document retrievable after a payload is inserted while constraining its
perplexity, followed by Context-Adaptive Trigger Generation (CATG), which
uses a frozen LLM to fuse a trigger naturally into the document's own
content. Under a one-poisoned-document-per-query evaluation on Natural
Questions, the paper reports high retrieval hit-rates (84.6\% / 81.3\%
HR@10 across two settings) together with substantial attack success at
near-benign perplexity. Its stated limitation is a white-box assumption
--- CBS needs gradient access to the retriever --- and evaluation on a
single dataset family. This project's ``stealth collusion'' attack
regime (\Cref{sec:attack-simulator}) is directly inspired by
SilentRetrieval's core idea --- camouflage indistinguishable from genuine
content --- adapted to a TF-IDF rather than a gradient-optimised setting,
since this project's environment has no route to a live embedding model
to attack with gradients.

\subsubsection{TriShieldRAG (arXiv:2607.23838, 2026)}
\label{sec:trishield-lit}

``TriShieldRAG: 3 Rings, One Blind Spot in Layered Defenses for
Retrieval-Augmented Generation'' proposes a three-ring, defense-in-depth
pipeline: an Ingest Guard that screens each candidate document with a
combined perplexity/repetition score, a boilerplate pattern-match score,
and an embedding-outlier score; a Retrieval Scorer that re-ranks the
surviving top-$k$ by a provenance-and-consistency-weighted trust score;
and a Cross-LLM Consensus stage that polls a panel of architecturally
diverse language models and allows one bounded re-retrieval on
disagreement. Evaluated against a non-adaptive PoisonedRAG-style attacker
over a 5,000-document knowledge base, the full pipeline reduces attack
success from roughly 91\% to roughly 13\%. The paper's most important
contribution for this project, however, is a negative result stated as
plainly as the positive one: agreement across the model panel is not the
same thing as correctness, since if every model in the panel is shown the
same poisoned retrieved context, a high agreement score and a high
attack-success rate can occur together. \Cref{sec:ring2} shows how this
project's own Ring~2 was redesigned specifically around that finding.

\subsection{Outcome of Literature Review}

Reading the six papers as a set, rather than in isolation, exposes a
consistent pattern: every one relies on a single primary signal, and
every one has a critical, paper-stated blind spot along a different
axis. No individual signal --- geometric anomaly, majority consensus,
single-document counterfactual influence, or cross-model agreement --- is
sufficient on its own.

\begin{table}[htbp]
\centering
\small
\begin{tabularx}{\textwidth}{@{}l Y Y@{}}
\toprule
\textbf{Paper} & \textbf{Primary Signal} & \textbf{Critical Blind Spot} \\
\midrule
DRS & Embedding-geometry anomaly at ingestion (low-variance-direction shift) & No query-time protection; the paper's own adaptive-attack experiment shows DRS-regularised poisoning reduces detection by roughly 15\% while keeping the attack effective \\
ShieldRAG (Push and Pull) & Query-time majority-consensus reshaping of the embedding space & Assumes benign information holds the majority; once poison documents form the majority the mechanism reshapes toward the poisoned consensus instead of away from it \\
RAGuard / ZKIP & Single-document leave-one-out counterfactual influence on the decoded answer & Blind to coordinated multi-document collusion --- removing one colluding document alone leaves the others' corroboration intact \\
PIDP-Attack & Dual-path attack combining a query-suffix trigger with corpus poisoning & Not a defense --- its significance is that a defense must protect the query path and the corpus path simultaneously, or the untouched path defeats it \\
SilentRetrieval & Fluency- and query-adaptive stealth poison (CBS + CATG) & Requires white-box access to the retriever's gradients to construct CBS triggers; defenses tuned only to catch obvious lexical or perplexity anomalies miss it by construction \\
TriShieldRAG & Three-ring layered defense (ingest guard, retrieval scorer, cross-model consensus) & ``False consensus'' --- the paper's own reported condition where cross-model agreement stays high while attack success also stays high, because every ring evaluates the same poisoned retrieved context \\
\bottomrule
\end{tabularx}
\caption{Six surveyed papers --- primary signal and critical blind spot.}
\label{tab:literature}
\end{table}

The key synthesis this project draws from \Cref{tab:literature} is that a
defense must combine independent, complementary signals at different
pipeline stages --- query, corpus, retrieval, and generation --- and must
be honest about the fact that some camouflaged attacks are, within a
single query's information content, structurally indistinguishable from
genuine agreement. That last point turns out, in \Cref{sec:results}, to
be more than a theoretical caveat: it is measured directly in this
project's own results.

\subsection{Problem Statement}

\textbf{Given} a RAG system whose retrieval corpus is vulnerable to (a)
standard single-document poisoning, (b) compound query-suffix-plus-corpus
poisoning (PIDP-style), (c) coordinated multi-document collusion, (d)
semantically fluent stealth collusion (SilentRetrieval-style camouflage
with no lexical tell), and (e) a single silently-camouflaged poison
document, \textbf{the problem} addressed by this project is to design,
implement, and empirically evaluate a defense framework that:

\begin{enumerate}
\item operates on real text with a real, fixed embedding space --- no
  synthetic noise vectors and no hand-tuned ``evasion bias'' parameters
  standing in for a detector's actual behaviour;
\item detects and measurably mitigates all five attack types listed
  above, individually and not only in aggregate;
\item calibrates every threshold from measured, held-out clean-data
  statistics rather than adjusting a parameter until a target headline
  number appears;
\item reports honest residual limitations --- including non-zero attack
  success where it genuinely occurs --- rather than presenting a
  suspiciously perfect result;
\item provides a per-component ablation showing which specific mechanism
  is responsible for which part of the final result; and
\item uses a multi-seed statistical methodology (independent corpora,
  confidence intervals) rather than a single run.
\end{enumerate}

\subsection{Research Objectives}
\label{sec:objectives}

\begin{enumerate}
\item[\textbf{O1}] --- Implement a four-ring defense pipeline
  (OmniGuard-RAG) in which each ring is a direct, traceable response to
  one specific gap identified in \Cref{sec:literature}.
\item[\textbf{O2}] --- Implement five baseline defenses (Vanilla RAG,
  DRS-only, ShieldRAG-only, RAGuard/ZKIP, TriShield) inside one common
  evaluation framework, so the comparison in \Cref{sec:results} is fair
  --- every system sees the same corpus, the same queries, and the same
  attacks.
\item[\textbf{O3}] --- Construct five realistic attack regimes (standard,
  PIDP, collusion, stealth collusion, silent) using query-targeted,
  factually-grounded poison generation rather than abstract placeholder
  text.
\item[\textbf{O4}] --- Conduct a multi-seed evaluation (8 independent
  seeds $\times$ 200 queries = 1,600 queries per system) reporting mean
  $\pm$ 95\% confidence interval for every metric.
\item[\textbf{O5}] --- Perform a per-ring ablation ladder that isolates
  each ring's individual contribution, deliberately excluding the
  persistent trust store so it cannot mask a ring's own limitation.
\item[\textbf{O6}] --- Report both an LLM-call-count estimate and a
  real, measured wall-clock orchestration latency for every system, since
  \Cref{sec:compute-cost} shows the two do not always agree.
\item[\textbf{O7}] --- Document every bug found and fixed during
  development, with its root cause and how it was measured, as part of
  the project's own methodology rather than as an appendix afterthought.
\end{enumerate}
